% Every chapter starts on a new page and begins with at least one paragraph
% before its first subchapter.
\chapter{Domain analysis}
\placeholder{Explain the domain, the existing approaches, the user needs and the requirements before describing the design. This opening paragraph is mandatory.}

\section{Problem and user needs}
\placeholder{Describe the people, their workflow and the concrete difficulty being addressed. At least one subchapter per page; two subchapters must not share a page.}

% Never start a page with a table or figure: introduce it with a short paragraph first.
\placeholder{According to Figure \ref{fig:context}, the system receives requests from the users and returns the processed result.}

\begin{figure}[htbp]
\centering
\includegraphics[width=0.45\textwidth]{example-diagram}
\caption{\placeholder{Figure title}}\label{fig:context}
\end{figure}

% Figures are numbered per chapter (1.1, 1.2, ...); the caption goes below the figure.
\placeholder{As shown in Table \ref{tab:requirements}, each requirement is linked to the user need that motivates it.}

% Tables: caption above, right aligned; table text is 11 pt with single spacing.
\begin{table}[htbp]
\caption{\placeholder{Table title}}\label{tab:requirements}
\centering
\begin{tabularx}{\textwidth}{|Y|Y|Y|Y|}\hline
\multicolumn{1}{|c|}{\textbf{A}} & \multicolumn{1}{c|}{\textbf{B}} & \multicolumn{1}{c|}{\textbf{C}} & \multicolumn{1}{c|}{\textbf{D}}\\\hline
 & & & \placeholder{Text 11 pt}\\\hline
 & & & \placeholder{Spacing 1}\\\hline
\end{tabularx}
\end{table}

\section{Existing approaches}
\placeholder{Compare relevant approaches using cited evidence \cite{example-article}.}

\section{Requirements}
\placeholder{List the functional and non-functional requirements that follow from the analysis.}
